\documentclass[11pt]{article}
\usepackage[a4paper,margin=2cm]{geometry}
\usepackage{amsmath,amssymb}
\usepackage[hidelinks]{hyperref}
\allowdisplaybreaks
\title{Qpu.General.Vocabulary}
\author{Tsvetan Rouschev}
\date{@uuidna/qpu v1.1.1, CC-BY-NC-ND-4.0}
\begin{document}
\maketitle
\noindent\url{https://github.com/uuidna/qpu/blob/v1.1.1/lean/Qpu/General/Vocabulary.lean}

\paragraph{\texttt{Vocabulary.adiabatic\_gap\_open}}
3.4.15 adiabatic quantum computing: the gap along the path from the first Hamiltonian to the last never closes.
\[ \forall (n\ k : \mathbb{N}), 0 < n \to k \le n \to 0 < \mathrm{Vocabulary.gapSquared}\ n\ k \]
\noindent axioms: none; proof: \url{https://github.com/uuidna/qpu/blob/v1.1.1/lean/Qpu/General/Vocabulary.lean\#L277}

\paragraph{\texttt{Vocabulary.basis\_length}}
3.1.4 simulator: a simulator of `n` qubits holds two to the power `n` amplitudes.
\[ \forall (n : \mathbb{N}), (\mathrm{Vocabulary.basis}\ n)\mathrm{.length} = 2 ^{n} \]
\noindent axioms: \texttt{propext}, \texttt{Quot.sound}; proof: \url{https://github.com/uuidna/qpu/blob/v1.1.1/lean/Qpu/General/Vocabulary.lean\#L33}

\paragraph{\texttt{Vocabulary.basis\_orthonormal}}
3.2.9 basis states: the basis states have unit length and are orthogonal.
\[ \mathrm{Vocabulary.inner}\ \mathrm{Vocabulary.e0}\ \mathrm{Vocabulary.e0} = 1 \land \mathrm{Vocabulary.inner}\ \mathrm{Vocabulary.e1}\ \mathrm{Vocabulary.e1} = 1 \land \mathrm{Vocabulary.inner}\ \mathrm{Vocabulary.e0}\ \mathrm{Vocabulary.e1} = 0 \]
\noindent axioms: none; proof: \url{https://github.com/uuidna/qpu/blob/v1.1.1/lean/Qpu/General/Vocabulary.lean\#L133}

\paragraph{\texttt{Vocabulary.bosons\_bunch}}
3.5.8 boson sampling: the permanent of the beam splitter is zero where its determinant is not, so two bosons never leave by different ports.
\[ \mathrm{Vocabulary.Hd.perm} = 0 \land \mathrm{Vocabulary.Hd.det} = -2 \]
\noindent axioms: none; proof: \url{https://github.com/uuidna/qpu/blob/v1.1.1/lean/Qpu/General/Vocabulary.lean\#L302}

\paragraph{\texttt{Vocabulary.cluster\_stabilizers\_commute}}
3.4.14 one-way quantum computing: the cluster state has one stabilizer for each qubit, and they all commute.
\[ \mathrm{Vocabulary.clusterStabilizers.length} = 3 \land (\mathrm{Vocabulary.clusterStabilizers.all}\ \lambda \ p \mapsto \mathrm{Vocabulary.clusterStabilizers.all}\ \lambda \ q \mapsto \mathrm{Vocabulary.commutes}\ p\ q) = \mathrm{true} \]
\noindent axioms: none; proof: \url{https://github.com/uuidna/qpu/blob/v1.1.1/lean/Qpu/General/Vocabulary.lean\#L269}

\paragraph{\texttt{Vocabulary.coding\_threshold}}
3.4.18 quantum coding, threshold: with a logical error of three times the square of the physical error, coding helps below one third and harms above it.
\[ 3 \cdot 100 \cdot 100 < 100 \cdot 1000 \land 400 \cdot 1000 < 3 \cdot 400 \cdot 400 \]
\noindent axioms: none; proof: \url{https://github.com/uuidna/qpu/blob/v1.1.1/lean/Qpu/General/Vocabulary.lean\#L285}

\paragraph{\texttt{Vocabulary.coherence\_bounded\_by\_relaxation}}
3.2.20 coherence time: twice the dephasing rate is at least the relaxation rate, so the coherence time is at most twice the relaxation time.
\[ \forall (\mathrm{relaxation}\ \mathrm{pure} : \mathbb{N}), \mathrm{relaxation} \le \mathrm{Vocabulary.twiceDephasing}\ \mathrm{relaxation}\ \mathrm{pure} \land (\mathrm{Vocabulary.twiceDephasing}\ \mathrm{relaxation}\ \mathrm{pure} = \mathrm{relaxation} \leftrightarrow \mathrm{pure} = 0) \]
\noindent axioms: \texttt{propext}, \texttt{Classical.choice}, \texttt{Quot.sound}; proof: \url{https://github.com/uuidna/qpu/blob/v1.1.1/lean/Qpu/General/Vocabulary.lean\#L174}

\paragraph{\texttt{Vocabulary.coherence\_interferes}}
3.2.18 quantum coherence: two Hadamard gates in a row return the basis state, because the second amplitude cancels.
\[ \mathrm{Vocabulary.Hd.apply} (\mathrm{Vocabulary.Hd.apply}\ \mathrm{Vocabulary.e0}) = (2, 0) \]
\noindent axioms: none; proof: \url{https://github.com/uuidna/qpu/blob/v1.1.1/lean/Qpu/General/Vocabulary.lean\#L158}

\paragraph{\texttt{Vocabulary.decoherence\_loses\_interference}}
3.2.19 decoherence: with the phase lost between the two gates the outcomes are even, where coherence gave one outcome only.
\[ \mathrm{Vocabulary.mixes} (\mathrm{Vocabulary.weights} (\mathrm{Vocabulary.Hd.apply}\ \mathrm{Vocabulary.e0})) = (2, 2) \land \mathrm{Vocabulary.weights} (\mathrm{Vocabulary.Hd.apply} (\mathrm{Vocabulary.Hd.apply}\ \mathrm{Vocabulary.e0})) = (4, 0) \]
\noindent axioms: none; proof: \url{https://github.com/uuidna/qpu/blob/v1.1.1/lean/Qpu/General/Vocabulary.lean\#L167}

\paragraph{\texttt{Vocabulary.eigenspace\_of\_identity}}
3.2.15 eigenspace: every state is kept by the identity, so its eigenspace for the value one is the whole space.
\[ \forall (v : \mathrm{Vocabulary.Vec}), \mathrm{Vocabulary.I.apply}\ v = v \]
\noindent axioms: \texttt{propext}; proof: \url{https://github.com/uuidna/qpu/blob/v1.1.1/lean/Qpu/General/Vocabulary.lean\#L151}

\paragraph{\texttt{Vocabulary.eigenstates\_of\_x}}
3.2.13 eigenstate: the two superpositions are the eigenstates of X, kept up to a sign.
\[ \mathrm{Vocabulary.X.apply} (1, 1) = (1, 1) \land \mathrm{Vocabulary.X.apply} (1, -1) = (-1, 1) \]
\noindent axioms: none; proof: \url{https://github.com/uuidna/qpu/blob/v1.1.1/lean/Qpu/General/Vocabulary.lean\#L145}

\paragraph{\texttt{Vocabulary.eigenvalues\_of\_z}}
3.2.14 eigenvalue: Z has the eigenvalues one and minus one, so its trace is zero and its determinant minus one.
\[ \mathrm{Vocabulary.Z.apply}\ \mathrm{Vocabulary.e0} = \mathrm{Vocabulary.e0} \land \mathrm{Vocabulary.Z.apply}\ \mathrm{Vocabulary.e1} = (0, -1) \land \mathrm{Vocabulary.Z.trace} = 0 \land \mathrm{Vocabulary.Z.det} = -1 \]
\noindent axioms: none; proof: \url{https://github.com/uuidna/qpu/blob/v1.1.1/lean/Qpu/General/Vocabulary.lean\#L148}

\paragraph{\texttt{Vocabulary.fidelity\_symmetric}}
3.3.11 fidelity: the overlap of two states does not depend on their order.
\[ \forall (u\ v : \mathrm{Vocabulary.Vec}), \mathrm{Vocabulary.inner}\ u\ v = \mathrm{Vocabulary.inner}\ v\ u \]
\noindent axioms: \texttt{propext}; proof: \url{https://github.com/uuidna/qpu/blob/v1.1.1/lean/Qpu/General/Vocabulary.lean\#L234}

\paragraph{\texttt{Vocabulary.gate\_fidelity}}
3.4.4 gate fidelity: the scaled Hadamard gate overlaps itself by four and a flip by two, so the flip is not that gate.
\[ \mathrm{Vocabulary.traceOverlap}\ \mathrm{Vocabulary.Hd}\ \mathrm{Vocabulary.Hd} = 4 \land \mathrm{Vocabulary.traceOverlap}\ \mathrm{Vocabulary.Hd}\ \mathrm{Vocabulary.X} = 2 \]
\noindent axioms: none; proof: \url{https://github.com/uuidna/qpu/blob/v1.1.1/lean/Qpu/General/Vocabulary.lean\#L246}

\paragraph{\texttt{Vocabulary.hamiltonian\_energies}}
3.2.12 Hamiltonian: the Hamiltonian gives the ground state energy minus one and the excited state energy one.
\[ \mathrm{Vocabulary.hamiltonian.apply}\ \mathrm{Vocabulary.e0} = (-1, 0) \land \mathrm{Vocabulary.hamiltonian.apply}\ \mathrm{Vocabulary.e1} = (0, 1) \]
\noindent axioms: none; proof: \url{https://github.com/uuidna/qpu/blob/v1.1.1/lean/Qpu/General/Vocabulary.lean\#L142}

\paragraph{\texttt{Vocabulary.key\_sifting}}
3.6.2 quantum cryptography: the receiver reads the sender's bit exactly when the two chose the same basis.
\[ \forall (\mathrm{bit}\ \mathrm{basis} : \mathrm{Bool}), \mathrm{Vocabulary.measure}\ \mathrm{basis} (\mathrm{Vocabulary.prepare}\ \mathrm{bit}\ \mathrm{basis}) = \mathrm{some}\ \mathrm{bit} \land \mathrm{Vocabulary.measure} (\neg \mathrm{basis}) (\mathrm{Vocabulary.prepare}\ \mathrm{bit}\ \mathrm{basis}) = \mathrm{none} \]
\noindent axioms: none; proof: \url{https://github.com/uuidna/qpu/blob/v1.1.1/lean/Qpu/General/Vocabulary.lean\#L313}

\paragraph{\texttt{Vocabulary.language\_evaluates}}
3.1.6 programming language: a term of the language has one meaning, and addition of terms is addition of meanings.
\[ \forall (a\ b : \mathrm{Vocabulary.Term}), (\mathrm{a.add}\ b)\mathrm{.eval} = \mathrm{a.eval} + \mathrm{b.eval} \land ((\mathrm{Vocabulary.Term.lit}\ 2)\mathrm{.add} ((\mathrm{Vocabulary.Term.lit}\ 3)\mathrm{.mul} (\mathrm{Vocabulary.Term.lit}\ 4)))\mathrm{.eval} = 14 \]
\noindent axioms: none; proof: \url{https://github.com/uuidna/qpu/blob/v1.1.1/lean/Qpu/General/Vocabulary.lean\#L51}

\paragraph{\texttt{Vocabulary.linear\_system\_solved}}
3.5.5 linear systems solving: the system has exactly one solution, and it is the one stated.
\[ \forall (v : \mathrm{Vocabulary.Vec}), \mathrm{Vocabulary.system.apply}\ v = (5, 10) \leftrightarrow v = (1, 3) \]
\noindent axioms: \texttt{propext}, \texttt{Quot.sound}; proof: \url{https://github.com/uuidna/qpu/blob/v1.1.1/lean/Qpu/General/Vocabulary.lean\#L293}

\paragraph{\texttt{Vocabulary.logical\_qudit\_corrects}}
3.3.8 logical qudit: a three-level digit repeated three times survives any one error, at any position and of any value.
\[ ((\mathrm{List.range}\ 3)\mathrm{.all}\ \lambda \ d \mapsto (\mathrm{List.range}\ 3)\mathrm{.all}\ \lambda \ \mathrm{position} \mapsto (\mathrm{List.range}\ 3)\mathrm{.all}\ \lambda \ \mathrm{value} \mapsto \mathrm{Vocabulary.majority3} (\mathrm{Vocabulary.corrupt} (\mathrm{Vocabulary.encode3}\ d) \mathrm{position}\ \mathrm{value}) = d) = \mathrm{true} \]
\noindent axioms: \texttt{propext}; proof: \url{https://github.com/uuidna/qpu/blob/v1.1.1/lean/Qpu/General/Vocabulary.lean\#L213}

\paragraph{\texttt{Vocabulary.logical\_x}}
3.3.9 logical operator: flipping every physical bit flips the encoded bit and nothing else.
\[ \forall (b : \mathrm{Bool}), \mathrm{Vocabulary.decodeBit} (\mathrm{List.map} (\lambda \ x \mapsto \neg x) (\mathrm{Vocabulary.encodeBit}\ b)) = \neg b \]
\noindent axioms: \texttt{propext}; proof: \url{https://github.com/uuidna/qpu/blob/v1.1.1/lean/Qpu/General/Vocabulary.lean\#L226}

\paragraph{\texttt{Vocabulary.mitigation\_extrapolates}}
3.4.7 quantum error mitigation: from a value measured at one and at two units of noise, the value without noise is twice the first minus the second.
\[ \forall (\mathrm{value}\ \mathrm{noise} : \mathbb{Z}), 2 \cdot (\mathrm{value} - \mathrm{noise}) - (\mathrm{value} - 2 \cdot \mathrm{noise}) = \mathrm{value} \]
\noindent axioms: \texttt{propext}, \texttt{Quot.sound}; proof: \url{https://github.com/uuidna/qpu/blob/v1.1.1/lean/Qpu/General/Vocabulary.lean\#L249}

\paragraph{\texttt{Vocabulary.model\_fits}}
3.1.1 model: the model with parameter three reproduces every observation.
\[ (\mathrm{Vocabulary.observations.all}\ \lambda \ x \mapsto \mathrm{match}\ x\ \mathrm{with} | (x, y) \mapsto \mathrm{Vocabulary.model}\ 3\ x = y) = \mathrm{true} \]
\noindent axioms: none; proof: \url{https://github.com/uuidna/qpu/blob/v1.1.1/lean/Qpu/General/Vocabulary.lean\#L21}

\paragraph{\texttt{Vocabulary.observable\_symmetric}}
3.2.11 quantum operator: the Pauli operators X and Z are their own transpose, so they are observables.
\[ \mathrm{Vocabulary.X.transpose} = \mathrm{Vocabulary.X} \land \mathrm{Vocabulary.Z.transpose} = \mathrm{Vocabulary.Z} \land \mathrm{Vocabulary.Hd.transpose} = \mathrm{Vocabulary.Hd} \]
\noindent axioms: none; proof: \url{https://github.com/uuidna/qpu/blob/v1.1.1/lean/Qpu/General/Vocabulary.lean\#L136}

\paragraph{\texttt{Vocabulary.operator\_linear}}
3.2.2 operator: an operator is linear, so it acts on a sum as the sum of its actions.
\[ \forall (m : \mathrm{Vocabulary.Mat}) (u\ v : \mathrm{Vocabulary.Vec}), \mathrm{m.apply} (\mathrm{u.1} + \mathrm{v.1}, \mathrm{u.2} + \mathrm{v.2}) = ((\mathrm{m.apply}\ u)\mathrm{.1} + (\mathrm{m.apply}\ v)\mathrm{.1}, (\mathrm{m.apply}\ u)\mathrm{.2} + (\mathrm{m.apply}\ v)\mathrm{.2}) \]
\noindent axioms: \texttt{propext}, \texttt{Quot.sound}; proof: \url{https://github.com/uuidna/qpu/blob/v1.1.1/lean/Qpu/General/Vocabulary.lean\#L116}

\paragraph{\texttt{Vocabulary.parameter\_unique}}
3.1.2 model parameter: the parameter that fits the first observation is unique.
\[ \forall (w : \mathbb{Z}), \mathrm{Vocabulary.model}\ w\ 1 = 3 \to w = 3 \]
\noindent axioms: \texttt{propext}; proof: \url{https://github.com/uuidna/qpu/blob/v1.1.1/lean/Qpu/General/Vocabulary.lean\#L24}

\paragraph{\texttt{Vocabulary.physical\_qudit\_levels}}
3.3.6 physical qudit: the transitions between neighbouring levels differ, so each of three levels is addressed alone.
\[ \mathrm{Vocabulary.level}\ 5000\ 200\ 1 - \mathrm{Vocabulary.level}\ 5000\ 200\ 0 = 5000 \land \mathrm{Vocabulary.level}\ 5000\ 200\ 2 - \mathrm{Vocabulary.level}\ 5000\ 200\ 1 = 4800 \land \mathrm{Vocabulary.level}\ 5000\ 200\ 3 - \mathrm{Vocabulary.level}\ 5000\ 200\ 2 = 4600 \]
\noindent axioms: none; proof: \url{https://github.com/uuidna/qpu/blob/v1.1.1/lean/Qpu/General/Vocabulary.lean\#L197}

\paragraph{\texttt{Vocabulary.process\_fidelity}}
3.3.12 process fidelity: a process overlaps itself by the dimension, and the identity does not overlap a flip.
\[ \mathrm{Vocabulary.traceOverlap}\ \mathrm{Vocabulary.I}\ \mathrm{Vocabulary.I} = 2 \land \mathrm{Vocabulary.traceOverlap}\ \mathrm{Vocabulary.X}\ \mathrm{Vocabulary.X} = 2 \land \mathrm{Vocabulary.traceOverlap}\ \mathrm{Vocabulary.I}\ \mathrm{Vocabulary.X} = 0 \]
\noindent axioms: none; proof: \url{https://github.com/uuidna/qpu/blob/v1.1.1/lean/Qpu/General/Vocabulary.lean\#L241}

\paragraph{\texttt{Vocabulary.projectors}}
3.2.17 projective measurement: each projector applied twice is itself, the two exclude each other, and they sum to the identity.
\[ \mathrm{Vocabulary.P0.mul}\ \mathrm{Vocabulary.P0} = \mathrm{Vocabulary.P0} \land \mathrm{Vocabulary.P1.mul}\ \mathrm{Vocabulary.P1} = \mathrm{Vocabulary.P1} \land \mathrm{Vocabulary.P0.mul}\ \mathrm{Vocabulary.P1} = \{ a := 0, b := 0, c := 0, d := 0 \} \land \mathrm{Vocabulary.P0.add}\ \mathrm{Vocabulary.P1} = \mathrm{Vocabulary.I} \]
\noindent axioms: none; proof: \url{https://github.com/uuidna/qpu/blob/v1.1.1/lean/Qpu/General/Vocabulary.lean\#L155}

\paragraph{\texttt{Vocabulary.quantum\_discrete}}
3.2.5 quantum, noun: levels that are whole multiples of a quantum have no level between two neighbours.
\[ \forall (q\ n\ m : \mathbb{N}), n \cdot q < m \cdot q \to (n + 1) \cdot q \le m \cdot q \]
\noindent axioms: \texttt{propext}, \texttt{Classical.choice}, \texttt{Quot.sound}; proof: \url{https://github.com/uuidna/qpu/blob/v1.1.1/lean/Qpu/General/Vocabulary.lean\#L122}

\paragraph{\texttt{Vocabulary.qudit\_register}}
3.3.4 qudit: a register of qudits has `d` to the power `n` basis states, and with two levels it is a register of qubits.
\[ \forall (n : \mathbb{N}), \mathrm{Vocabulary.registerDim}\ 2\ n = (\mathrm{Vocabulary.basis}\ n)\mathrm{.length} \land \mathrm{Vocabulary.registerDim}\ 3\ 2 = 9 \]
\noindent axioms: \texttt{propext}, \texttt{Quot.sound}; proof: \url{https://github.com/uuidna/qpu/blob/v1.1.1/lean/Qpu/General/Vocabulary.lean\#L190}

\paragraph{\texttt{Vocabulary.relaxation\_halves}}
3.2.21 relaxation time: a population halves with each half-life and is gone after finitely many.
\[ \mathrm{Vocabulary.remaining}\ 1024\ 1 = 512 \land \mathrm{Vocabulary.remaining}\ 1024\ 10 = 1 \land \mathrm{Vocabulary.remaining}\ 1024\ 11 = 0 \]
\noindent axioms: none; proof: \url{https://github.com/uuidna/qpu/blob/v1.1.1/lean/Qpu/General/Vocabulary.lean\#L182}

\paragraph{\texttt{Vocabulary.sensing\_gain}}
3.6.4 quantum sensing: the precision `n` entangled probes reach is at least the precision `n` separate probes reach.
\[ \forall (n : \mathbb{N}), n \le n \cdot n \]
\noindent axioms: \texttt{propext}; proof: \url{https://github.com/uuidna/qpu/blob/v1.1.1/lean/Qpu/General/Vocabulary.lean\#L318}

\paragraph{\texttt{Vocabulary.state\_fidelity}}
3.3.10 state fidelity: a basis state overlaps itself fully and the other basis state not at all.
\[ \mathrm{Vocabulary.inner}\ \mathrm{Vocabulary.e0}\ \mathrm{Vocabulary.e0} = 1 \land \mathrm{Vocabulary.inner}\ \mathrm{Vocabulary.e0}\ \mathrm{Vocabulary.e1} = 0 \land \mathrm{Vocabulary.inner} (1, 1) \mathrm{Vocabulary.e0} \cdot \mathrm{Vocabulary.inner} (1, 1) \mathrm{Vocabulary.e0} \cdot 2 = \mathrm{Vocabulary.inner} (1, 1) (1, 1) \]
\noindent axioms: none; proof: \url{https://github.com/uuidna/qpu/blob/v1.1.1/lean/Qpu/General/Vocabulary.lean\#L230}

\paragraph{\texttt{Vocabulary.superposition}}
3.2.8 quantum superposition: the Hadamard gate takes a basis state to a state with both amplitudes, equal.
\[ \mathrm{Vocabulary.Hd.apply}\ \mathrm{Vocabulary.e0} = (1, 1) \land \mathrm{Vocabulary.Hd.apply}\ \mathrm{Vocabulary.e1} = (1, -1) \]
\noindent axioms: none; proof: \url{https://github.com/uuidna/qpu/blob/v1.1.1/lean/Qpu/General/Vocabulary.lean\#L130}

\paragraph{\texttt{Vocabulary.system\_dimension}}
3.2.6 quantum system: two systems taken as one have the product of their dimensions.
\[ \forall (m\ n : \mathbb{N}), (\mathrm{Vocabulary.basis} (m + n))\mathrm{.length} = (\mathrm{Vocabulary.basis}\ m)\mathrm{.length} \cdot (\mathrm{Vocabulary.basis}\ n)\mathrm{.length} \]
\noindent axioms: \texttt{propext}, \texttt{Quot.sound}; proof: \url{https://github.com/uuidna/qpu/blob/v1.1.1/lean/Qpu/General/Vocabulary.lean\#L126}
\end{document}
